\item \textbf{ACTIVIDAD} Considere los 10 objetos de nuestro sistema solar en la lista más abajo. Utilizando las temperaturas de la superficie enumeradas para estos objetos, determine la velocidad promedio de los átomos de helio ubicados en las atmósferas de estos objetos. Usando datos planetarios, determine la velocidad de escape para cada uno de los objetos. Finalmente, tome la relación de la velocidad de escape a la velocidad promedio de los átomos de helio para cada objeto. (a) ¿Cuál es el valor típico de esta relación para un objeto con poca o ninguna atmósfera? (b) ¿Cuál es el valor típico de esta relación para que un objeto tenga una atmósfera robusta, pero con poco o nada de helio presente? (c) ¿Cuál es el valor típico de esta relación para que un objeto tenga una atmósfera robusta con una cantidad significativa de helio?

\begin{center}
\begin{tabular}{|l|c|l|}
\hline
\textbf{Objeto} & \textbf{Temperatura de la superficie} & \textbf{Atmósfera} \\ 
 & \textbf{(del Instituto Lunar y Planetario)} & \\ \hline
Mercurio & 440 & Poca o nada \\
Venus & 741 & Robusta, poco He \\
Tierra & 288 & Robusta, poco He \\
Marte & 244 & Robusta, poco He \\
Ceres & 173 & Poca o nada \\
Júpiter & 165 & Robusta, mucho He \\
Saturno & 134 & Robusta, mucho He \\
Urano & 77 & Robusta, mucho He \\
Neptuno & 70 & Robusta, mucho He \\
Plutón & 40 & Poca o nada \\ \hline
\end{tabular}
\end{center}
